

\usepackage[utf8]{inputenc}
\usepackage[T1]{fontenc}

\usepackage{etoolbox}
\usepackage[en-US]{datetime2}

\usepackage{amsmath} %equations and spacings
\usepackage{amssymb}
\usepackage{mathcommand} % added by M.E.
\usepackage{amsfonts} %mathbb
\usepackage{amsthm} % theorems and proofs
\usepackage{mathtools} % theorems and proofs
\usepackage{multicol}
\usepackage{empheq}
\usepackage{appendix}


% \usepackage{easyReview}
% \usepackage{tgpagella}
% \usepackage{tgpagella,eulervm}
% \usepackage{newpxtext,newpxmath}
% \usepackage{newpxtext,eulervm}
% Custom commands
% Custom todo command, switch definitions if you want to quickly 
% \newcommand{\todo}[1]{{\color{red}\textbf{(#1)}}}
\newcommand{\todo}[1]{}


% Abbrieviations, taken from https://stackoverflow.com/questions/3282319/correct-way-to-define-macros-etc-ie-in-latex
% These definitions are used in CVPR's style package
\usepackage{xspace}

% Add a period to the end of an abbreviation unless there's one
% already, then \xspace.
\makeatletter
\DeclareRobustCommand\onedot{\futurelet\@let@token\@onedot}
\def\@onedot{\ifx\@let@token.\else.\null\fi\xspace}

\def\eg{{e.g}\onedot} 
\def\Eg{{E.g}\onedot}
\def\ie{{i.e}\onedot} 
\def\Ie{{I.e}\onedot}
\def\cf{{c.f}\onedot} 
\def\Cf{{C.f}\onedot}
\def\etc{{etc}\onedot} 
\def\vs{{vs}\onedot}
\def\wrt{{w.r.t}\onedot} 
\def\dof{{d.o.f}\onedot}
\def\etal{{et al}\onedot}
\makeatother



%%%%%%%%%%%%%%%%%%%%%%%%%%%%%%%%%%%%%%
%%% Definitions and lengths %%%
%%%%%%%%%%%%%%%%%%%%%%%%%%%%%%%%%%%%%%

% Font size (gets local font size)
\makeatletter
\newcommand*\myfontsize{\dimexpr\f@size pt\relax}
\makeatother

% British spellcheck software tends to complain about this
\newcommand{\centring}{\centering}
\newenvironment{itemise}{\begin{itemize}}{\end{itemize}}%
\usepackage{xcolor}
\usepackage{color}

%%%%%%%%%%%%%% MATLAB COLORS %%%%%%%%%%%%%%
\definecolor{softred}{HTML}{b7312c}
\definecolor{softyellow}{HTML}{f2a900}
\definecolor{softgreen}{HTML}{48a23f}
\definecolor{softblue}{HTML}{00a9e0}
\definecolor{softpurple}{HTML}{715091}
\definecolor{softgray}{HTML}{636569}

%%%%%%%%%%%%%% MATPLOTLIB COLORS %%%%%%%%%%%%%%
\definecolor{Tblue}{HTML}{1f77b4}
\definecolor{Torange}{HTML}{ff7f0e}
\definecolor{Tgreen}{HTML}{2ca02c}
\definecolor{Tred}{HTML}{d62728}
\definecolor{Tpurple}{HTML}{9467bd}
\definecolor{Tbrown}{HTML}{8c564b}

%%%%%%%%%%%%%% MISC COLORS %%%%%%%%%%%%%%
\definecolor{AdilsBlue}{HTML}{03a9f4}

%%some other colors

\definecolor{hlpink}{HTML}{FFACBD}
\definecolor{hlorange}{HTML}{FCD6B4}
\definecolor{hlyellow}{HTML}{FBF3C0}
\definecolor{hlwhite}{HTML}{FFFFFF}
\definecolor{hlblue}{HTML}{D3F4F3}
\definecolor{hlgreen}{HTML}{B1D29B}




% Make a custom legend for my plots
% argument #1: any options
\makeatletter
\newenvironment{customlegend}[1][]{%
    \begingroup
    % inits/clears the lists (which might be populated from previous
    % axes):
    \pgfplots@init@cleared@structures
    \pgfplotsset{#1}%
}{%
    % draws the legend:
    \pgfplots@createlegend
    \endgroup
}%

% makes \addlegendimage available (typically only available within an
% axis environment):
\def\addlegendimage{\pgfplots@addlegendimage}
\makeatother

% \todo{Remove this and scale the plots properly in tikz}
% Used for figure scaling in chapter 4.2
\newcommand{\tikzFigureScale}{0.55}


%% General convenience %%


%% Drawing
% \newcommand{\tikzdot}[2][red,fill=red]{\tikz[baseline=-0.5ex]\draw[#1,radius=#2] (0,0) circle ;}%
\newcommand{\drawbar}[1]{\begin{tikzpicture}\draw [white] (0,0) rectangle (\myfontsize, \myfontsize); \draw [#1] (0, 0.25*\myfontsize) -- (0.333*\myfontsize, 0.25*\myfontsize); \end{tikzpicture}}

\newcommand{\drawdash}[1]{\begin{tikzpicture}\draw [white] (0,0) rectangle (\myfontsize, \myfontsize);\draw [#1] (0, 0.25*\myfontsize) -- (0.333*\myfontsize, 0.25*\myfontsize);\draw [#1] (0.666*\myfontsize, 0.25*\myfontsize) -- (\myfontsize, 0.25*\myfontsize);\end{tikzpicture}}








\makeatletter
\renewenvironment{proof}[1][\relax]{\par
  \pushQED{\qed}%
  \normalfont \topsep6\p@\@plus6\p@\relax
  \trivlist
  \item[\hskip\labelsep\bfseries
    \ifx#1\relax \proofname\else#1\fi\@addpunct{.}]\ignorespaces
}{%
  \popQED\endtrivlist\@endpefalse
}
\makeatother

\DTMlangsetup{showdayofmonth=false}

\renewcommand{\qedsymbol}{$\blacksquare$}
% \usepackage{utils/texnax}
%% Nicely format and linebreak URLs in the bibliography (and elsewhere).
\usepackage{url}
%% Define a new 'leo' style for the package that will use a smaller font.
\makeatletter
\def\url@leostyle{%
  \@ifundefined{selectfont}{\def\UrlFont{\sf}}{\def\UrlFont{\small\ttfamily}}}
\makeatother
%% Now actually use the newly defined style.
\urlstyle{leo}


\usepackage{nicefrac} % nice inline fractions
\usepackage{bm}
\usepackage[version=4]{mhchem} % Chemistry
\usepackage{siunitx}
\usepackage{pifont} % used for \ding
\usepackage{lipsum} % acces to random text
% \usepackage{blindtext} % acces to random text

\usepackage{epigraph,varwidth}
\renewcommand{\epigraphsize}{\small}
\setlength{\epigraphwidth}{0.9\textwidth}
\renewcommand{\textflush}{center}
\renewcommand{\sourceflush}{center}
\renewcommand{\epigraphflush}{center}
% A useful addition
\newcommand{\epitextfont}{\itshape}
\newcommand{\episourcefont}{\scshape}

\makeatletter
\newsavebox{\epi@textbox}
\newsavebox{\epi@sourcebox}
\newlength\epi@finalwidth
\renewcommand{\epigraph}[2]{%
  \vspace{\beforeepigraphskip}
  {\epigraphsize\begin{\epigraphflush}
   \epi@finalwidth=\z@
   \sbox\epi@textbox{%
     \varwidth{\epigraphwidth}
     \begin{\textflush}\epitextfont#1\end{\textflush}
     \endvarwidth
   }%
   \epi@finalwidth=\wd\epi@textbox
   \sbox\epi@sourcebox{%
     \varwidth{\epigraphwidth}
     \begin{\sourceflush}\episourcefont#2\end{\sourceflush}%
     \endvarwidth
   }%
   \ifdim\wd\epi@sourcebox>\epi@finalwidth 
     \epi@finalwidth=\wd\epi@sourcebox
   \fi
   \leavevmode\vbox{
     \hb@xt@\epi@finalwidth{\hfil\box\epi@textbox}
     \vskip1.75ex
     \hrule height \epigraphrule
     \vskip.75ex
     \hb@xt@\epi@finalwidth{\hfil\box\epi@sourcebox}
   }%
   \end{\epigraphflush}
   \vspace{\afterepigraphskip}}}
\makeatother

%graphics
\usepackage{graphicx}
\usepackage{svg}
% \graphicspath{.}
%\usepackage[demo]{graphicx}

% These settings control how floats (figures) can occupy the page
% Taken from here: https://robjhyndman.com/hyndsight/latex-floats/
\setcounter{topnumber}{2} % max number of figs at top
\setcounter{bottomnumber}{2} % % max number of figs at bottom
\setcounter{totalnumber}{4} % max number of figs on 1 page
\renewcommand{\topfraction}{0.85} % max proportion of top figures can take
\renewcommand{\bottomfraction}{0.85} % max proportion of top figures can take
\renewcommand{\textfraction}{0.15} % min proportion of text on 1 page
\renewcommand{\floatpagefraction}{0.7} % mim fraction of floatpage that can have floats

%caption styles
\usepackage[small,bf]{caption}
\usepackage[labelformat=simple]{subcaption}
\renewcommand\thesubfigure{(\alph{subfigure})} % see subcaption doc

%% We want UK-english hyphenation patterns
\usepackage[british]{babel}
\usepackage{csquotes}

%Remove Paragraph indenting
\setlength{\parindent}{0.0in} 
\setlength{\parskip}{0.1in}
\setlength{\headheight}{21.75pt}

%% Use scalable, PostScript Type 1 versions of the Computer Modern fonts.
\usepackage{type1cm}
%% Replace the standard Computer Modern Typewriter font LaTeX uses
%% for monospace text with the PostScript font Adobe Courier.
\usepackage{courier}
\usepackage{ae,aecompl}
\usepackage{times}

% %% Redefine the font used for the section headings to
% %% Helvetica-Narrow Bold.
% \usepackage{sectsty}
% \allsectionsfont{\usefont{OT1}{phv}{bc}{n}\selectfont}

%to control the title formating and spaces around the titles
%spaces under sections
\usepackage[compact]{titlesec}
\titlespacing{\section}{0pt}{*0}{*0}
\titlespacing{\subsection}{0pt}{*0}{*0}
\titlespacing{\subsubsection}{0pt}{*0}{*0}

%Abbreviations or Acronyms
\usepackage{nomencl}
\renewcommand{\nomname}{List of Abbreviations}
\makenomenclature

% \usepackage[numbers]{cite}
\usepackage[numbers]{natbib}
\def\chosenbibstyle{IEEEtran}
\def\bibfilename{bibliography}
% \usepackage[style=numeric,citestyle=numeric,natbib=true,sorting=none]{biblatex}
% % Remove the following fields
% \DeclareSourcemap{
%   \maps[datatype=bibtex, overwrite]{
%     \map{
%       \step[fieldset=address, null]
%       \step[fieldset=location, null]
%       \step[fieldset=editor, null]
%     }
%   }
% }

% \addbibresource{zotero_references.bib} 
% \addbibresource{additional_references.bib} 

% % Turn off doi links in the bib
% \ExecuteBibliographyOptions{doi=false,url=false}
% % Make the titles a link
% \newbibmacro{string+doi}[1]{%
%   \iffieldundef{doi}
%   {\iffieldundef{url}{#1}{\href{\thefield{url}}{#1}}}
%   {\href{http://dx.doi.org/\thefield{doi}}{#1}}}
% \newbibmacro{eprint-nolink}[1]{%
%   \iffieldundef{eprint}
%   {\iffieldundef{url}{#1}{\href{\thefield{url}}{#1}}}
%   {\href{http://dx.doi.org/\thefield{doi}}{#1}}}
% \DeclareFieldFormat[misc,article,inproceedings,book,thesis,incollection,report]{title}{\usebibmacro{string+doi}{\mkbibquote{#1}}}


% % Only show arxiv eprints, but with no link
% \DeclareFieldFormat{eprint}{} 
% \DeclareFieldFormat{eprint:arxiv}{arXiv: \thefield{eprint}} 

% % ISSN is disabled via isbn=false, but I want isbn for books.
% \DeclareFieldFormat{issn}{} 

% Don't need to know the address of whatever
% \DeclareFieldFormat{location}{} 
% \DeclareFieldFormat{address}{} 

\usepackage{hyperref}
%\hypersetup{colorlinks=true,citecolor=blue}
\hypersetup{colorlinks=true,citecolor=Tblue,linkcolor=black,urlcolor=Tblue}
%To change box color around the links and citations, you have these other options:
%\hypersetup{citebordercolor=Violet,filebordercolor=Red,linkbordercolor=Blue}

% Cleveref must be loaded after hyperref
\usepackage[capitalise]{cleveref} % Auto references

\newcommand*\makeAlph[1]{\symbol{\numexpr64+#1}} % A,B,C,...


% Add Assumption type to cleveref 
\crefformat{assumption}{Assumption~#2#1#3}
\crefrangeformat{assumption}{Assumptions~(#3#1#4) to~(#5#2#6)}
\crefmultiformat{assumption}{Assumptions~#2#1#3}%
{ and~#2#1#3}{, #2#1#3}{ and~#2#1#3}



% Add Paper type to cleveref
%% Define a counter that we can use for references
\newcounter{paper}
%% Convenience command to label papers (prints a ref to the paper)
\newcommand{\labelpaper}[1]{{\refstepcounter{paper}\label{#1}}\cref{#1}}
%% Convenience command to turn numbers into capital letters (1->A, 2->B, ...)
%% Define how the reference will look (Paper A, Paper B, etc)
\crefformat{paper}{Paper~#2\makeAlph{#1}#3}
\crefrangeformat{paper}{Papers~#3\makeAlph{#1}#4 to~#5\makeAlph{#2}#6}
\crefmultiformat{paper}{Papers~#2\makeAlph{#1}#3}%
{ and~#2\makeAlph{#1}#3}{, #2\makeAlph{#1}#3}{ and~#2\makeAlph{#1}#3}



% Add RQ type to cleveref
%% Define a counter that we can use for references
\newcounter{rq}
% \renewcommand{\therq}{\Roman{rq}}
%% Convenience command to label papers (prints a ref to the paper)
\newcommand{\labelrq}[1]{{\refstepcounter{rq}\label{#1}}\cref{#1}}
%% Convenience command to turn numbers into capital letters (1->A, 2->B, ...)
%% Define how the reference will look (Paper A, Paper B, etc)
\crefformat{rq}{RQ~#2#1#3}
\crefrangeformat{rq}{RQs~#3#1#4 to~#5\makeRoman{#2}#6}
\crefmultiformat{rq}{RQs~#2#1#3}%
{ and~#2#1#3}{, #2#1#3}{ and~#2#1#3}

  
% Extra commands for tables, like \midrule
\usepackage{booktabs}

% Special table stuff
\usepackage{array}
\usepackage{multirow}
\usepackage{dcolumn}
\newcommand{\PreserveBackslash}[1]{\let\temp=\\#1\let\\=\temp}
\newcolumntype{d}{D{.}{.}{-2}}
\newcolumntype{C}[1]{>{\PreserveBackslash\centring}p{#1}}
\newcolumntype{R}[1]{>{\PreserveBackslash\raggedleft}p{#1}}
\newcolumntype{L}[1]{>{\PreserveBackslash\raggedright}p{#1}} 
\newcolumntype{x}[1]{>{\centring\let\newline\\\arraybackslash\hspace{0pt}}p{#1}} % You can specify cell width, but is centred


% Fancy chapter options, no options seems to reset the font set by sectsty
\usepackage{fncychap}




%graphics
\usepackage{epstopdf} %support for eps.
\DeclareGraphicsExtensions{.eps,.ps,.pdf,.png,.jpg}
\usepackage{float} % figure placing [H]


%Section Numbering and TOC depth
\setcounter{secnumdepth}{2}
\setcounter{tocdepth}{2}


%Algorithms
% \usepackage[ruled,vlined]{algorithm2e}  % Algorithms with style
\usepackage{algorithm}
\usepackage{algpseudocode}

% Number the algos by chapter
% algorithm2e has options for this, but isn't compatible with 
% algpseudocode. Oh well!
% https://tex.stackexchange.com/questions/124902/algorithm-with-chapter-number
\makeatletter 
\renewcommand\thealgorithm{\thechapter.\arabic{algorithm}} 
\@addtoreset{algorithm}{chapter} 
\makeatother


% Glossary
\usepackage[acronym, nomain, automake]{glossaries-extra} 
\renewcommand*{\glstextformat}[1]{\textcolor{black}{#1}}
\glssetcategoryattribute{acronym}{glossdesc}{title}
\setabbreviationstyle[acronym]{long-short}
% extra acronym stuff
\glsaddkey*{adjective}{???}{\glsentryadj}{\Glsentryadj}{\glsadjx}{\Glsadjx}{\GLSadjx}

\newcommand{\ADJglspostlinkhook}{\ifglsused{\glslabel}{}{ (\glsentryshort{\glslabel})}\glsunset{\glslabel}}
\newcommand{\glsadj}[1]{{\let\glspostlinkhook\ADJglspostlinkhook\ifglsused{#1}{\gls{#1}}{\glsadjx{#1}}}}
\newcommand{\Glsadj}[1]{{\let\glspostlinkhook\ADJglspostlinkhook\ifglsused{#1}{\Gls{#1}}{\Glsadjx{#1}}}}
\newcommand{\GLSadj}[1]{{\let\glspostlinkhook\ADJglspostlinkhook\ifglsused{#1}{\GLS{#1}}{\GLSadjx{#1}}}}
\newacronym[adjective={uniformly stable}]{us}{US}{uniform stability}
\newacronym[adjective={globally stable}]{gs}{GS}{global stability}
\newacronym[adjective={uniformly globally stable}]{ugs}{UGS}{uniform global stability}

\newacronym[adjective={asymptotically stable}]{as}{AS}{asymptotic stability}
\newacronym[adjective={uniformly asymptotically stable}]{uas}{UAS}{uniform asymptotic stability}
\newacronym[adjective={globally asymptotically stable}]{gas}{GAS}{global asymptotic stability}
\newacronym[adjective={uniformly globally asymptotically stable}]{ugas}{UGAS}{uniform global asymptotic stability}
\newacronym[adjective={semiglobally asymptotically stable}]{sgas}{SGAS}{semiglobal asymptotic stability}
\newacronym[adjective={uniformly semiglobally asymptotically stable}]{usgas}{USGAS}{uniform semiglobal asymptotic stability}
\newacronym[adjective={uniformly practically asymptotically stable}]{upas}{UPAS}{uniform pracitcal asymptotic stability}

\newacronym[adjective={exponentially stable}]{es}{ES}{exponential stability}
\newacronym[adjective={locally exponentially stable}]{les}{LES}{local exponential stability}
\newacronym[adjective={uniformly exponentially stable}]{ues}{UES}{uniform exponential stability}
\newacronym[adjective={globally exponentially stable}]{ges}{GES}{global exponential stability}
\newacronym[adjective={uniformly globally exponentially stable}]{uges}{UGES}{uniform global exponential stability}
\newacronym[adjective={semiglobally exponentially stable}]{sges}{SGES}{semiglobal exponential stability}
\newacronym[adjective={uniformly semiglobally exponentially stable}]{usges}{USGES}{uniform semiglobal exponential stability}
\newacronym[adjective={exponentially orbitally stable}]{eos}{EOS}{exponential orbital stability}

\newacronym[adjective={globally bounded}]{gb}{GB}{global boundedness}
\newacronym[adjective={ultimately bounded}]{ub}{UB}{ultimate boundedness}
\newacronym[adjective={uniformly ultimately bounded}]{uub}{UUB}{uniform ultimate boundedness}
\newacronym[adjective={uniformly globally bounded}]{ugb}{UGB}{uniform global boundedness}
\newacronym[adjective={uniformly globally ultimately bounded}]{ugub}{UGUB}{uniform global ultimate boundedness}

\newacronym[adjective={globally attractive}]{ga}{GA}{global attractivity}
\newacronym[adjective={semiglobally attractive}]{sga}{SGA}{semiglobal attractivity}

% \newacronym[adjective={practically asymptotically stable}]{pas}{PAS}{practical asymptotic stability}
% \newacronym[adjective={globally practically asymptotically stable}]{gpas}{GPAS}{global practical asymptotic stability}
% \newacronym[adjective={uniformly globally practically asymptotically stable}]{ugpas}{UGPAS}{uniform global practical asymptotic stability}
% \newacronym[adjective={semiglobally practically asymptotically stable}]{sgpas}{SGPAS}{semiglobal practical asymptotic stability}
% \newacronym[adjective={uniformly semiglobally practically asymptotically stable}]{uspas}{USPAS}{uniform semiglobal practical asymptotic stability}

\newacronym{ltv}{LTV}{linear time-varying}
\newacronym{lti}{LTI}{linear time-invariant}

\newacronym{clf}{CLF}{control-Lyapunov function}
\newacronym{lfc}{LFC}{Lyapunov function candidate}

\newacronym[adjective={persistently exciting}]{pe}{PE}{persistency of excitation}

% Control methods

\newacronym{mpc}{MPC}{model predictive control}
\newacronym{nmpc}{NMPC}{nonlinear model predictive control}
\newacronym{smc}{SMC}{sliding mode control}
\newacronym{stsmc}{STSMC}{super-twisting sliding mode control}
\newacronym{gsta}{GSTA}{generalized super-twisting algorithm}
\newacronym{rl}{RL}{reinforcement learning}
\newacronym{esc}{ESC}{extremum-seeking control}


% Snake/marine lingo
\newacronym{usv}{USV}{uncrewed surface vehicle}
\newacronym[shortplural={AUVs},longplural={autonomous underwater vehicles}]{auv}{AUV}{autonomous underwater vehicle}
\newacronym{rov}{ROV}{remotely operated vehicle}
\newacronym{vms}{VMS}{vehicle-manipulator system}

\newacronym[shortplural={USRs},longplural={underwater snake robots}]{usr}{USR}{underwater snake robot}
\newacronym{aiauv}{AIAUV}{articulated intervention autonomous underwater vehicle}
\newacronym{wec}{WEC}{wave energy conversion}

% Some other specific things

% Some general stuff
\newacronym{rms}{RMS}{root mean square}
\newacronym{ad}{AD}{automatic differentiation}
\newacronym{nlp}{NLP}{nonlinear programming}
\newacronym[longplural={linear independence constraint qualifications}]{licq}{LICQ}{linear independence constraint qualification}
\newacronym{nloc}{NLOC}{nonlinear optimal control}
\newacronym{lq}{LQ}{linear-quadratic}
\newacronym{qp}{QP}{quadratic programming}
\newacronym[longplural={bilinear matrix inequalities}]{bmi}{BMI}{bilinear matrix inequality}
\newacronym[longplural={linear matrix inequalities}]{lmi}{LMI}{linear matrix inequality}
\newacronym{ode}{ODE}{ordinary differential equation}
\newacronym{mpp}{MPP}{Moore-Penrose pseudoinverse}
\newacronym{mbs}{MBS}{multi-body system}
\newacronym{vpm}{VPM}{vortex particle-mesh}
\newacronym{cfd}{CFD}{computational fluid dynamics}
\newacronym[shortplural={VHCs},longplural={virtual holonomic constraints}]{vhc}{VHC}{virtual holonomic constraints}
\newacronym[shortplural={CPGs},longplural={central pattern generators}]{cpg}{CPG}{central pattern generator}
\newacronym{fsi}{FSI}{fluid-structure interaction}
\newacronym{wopc}{WOPC}{weather optimal position control}
\newacronym{wohc}{WOHC}{weather optimal heading control}
\newacronym{cm}{CM}{center of mass}
\newacronym{lcfl}{LCFL}{Lagrangian Courant-Friedrich-Levy}


\newacronym[adjective={uniformly globally attractive}]{uga}{UGA}{uniform global attractivity}
\newacronym[adjective={uniformly almost globally attractive}]{uaga}{UAGA}{uniform almost global attractivity}
\newacronym[adjective={globally pre-attractive}]{gpa}{GpA}{global pre-attractivity}
\newacronym[adjective={almost globally pre-attractive}]{agpa}{AGpA}{almost global pre-attractivity}
\newacronym[adjective={almost globally asymptotically stable}]{agas}{AGAS}{almost global asymptotic stability}
\newacronym[adjective={almost globally pre-asymptotically stable}]{agpas}{AGpAS}{almost global pre-asymptotic stability}
\newacronym[adjective={almost globally attractive}]{aga}{AGA}{almost global attractivity}
\newacronym[adjective={uniformly globally pre-asymptotically stable}]{ugpas}{UGpAS}{uniform global pre-asymptotic stability}
\newacronym[adjective={uniformly almost globally pre-asymptotically stable}]{uagpas}{UAGpAS}{uniform almost global pre-asymptotic stability}
\newacronym[adjective={uniformly almost globally asymptotically stable}]{uagas}{UAGAS}{uniform almost global asymptotic stability}


\makeglossaries

%%%%%%%%%%%%%%%%%%%%%%%%%%%%%%%%%%%%%%
%%% Math commands %%%
%%%%%%%%%%%%%%%%%%%%%%%%%%%%%%%%%%%%%%
%% Generally useful

% Vector or matrix (boldface)
% \newcommand\vm[1]{\boldsymbol{\mathrm{#1}}}
\newcommand\vm[1]{\bm{#1}}


% R^n
\newcommand\R[1][\empty]{\mathbb{R}^{#1}}

% C^n
\newcommand\C[1][\empty]{\mathbb{C}^{#1}}

% Identity
\newcommand\I[1][\empty]{\vm{I}_{#1}}

% Transpose
\gdef\T{^{\top}}

% Inverse
\gdef\inv{^{-1}}

% Moore-Penrose Pseudoinverse
\gdef\mpp{^{+}}

% Set of elements
\newcommand\set[1]{\left\{#1\right\}}

% Span of vectors (\span was taken, sadly :( )
\newcommand\spn[1]{\textnormal{span}\left(#1\right)}

% Range of matrix
\newcommand\ran[1]{\textnormal{range}\left(#1\right)}

% Orthogonal complement
\gdef\orth{^{\perp}}

% Better trigonometric functions if one wants


% \let\oldcos\cos
% \renewcommand\cos[1]{\oldcos{\left(#1\right)}}
% \newcommand\scos[1]{\oldcos{^2\left(#1\right)}}

% \let\oldsin\sin
% \renewcommand\sin[1]{\oldsin{\left(#1\right)}}
% \newcommand\ssin[1]{\oldsin{^2\left(#1\right)}}

% \let\oldtan\tan
% \renewcommand\tan[1]{\oldtan{\left(#1\right)}}
% \newcommand\stan[1]{\oldtan{^2\left(#1\right)}}

% Hadamard product (slightly unnecessary? less readable? more readable?)
\let\hprod\odot

% Hadamard power
\newcommand\hpow[1]{^{\circ#1}}

% Kronecker product
\let\kron\otimes

% Slightly more meaningful names for pos.def, neg.def notation

% Pos. def, "matrix greater than"
\let\mgt\succ

% Neg. def
\let\mlt\prec

% Pos. semidef.
\let\mgeq\succcurlyeq

% Neg. semidef
\let\mleq\preccurlyeq

% Better real and imaginary part if one wants

% \let\oldRe\Re
% \renewcommand\Re[1]{\oldRe\left(#1\right)}

% \let\oldIm\Im
% \renewcommand\Im[1]{\oldIm\left(#1\right)}

%% Slightly more specific useful stuff

% Vectors with optional index and derivatives
\renewmathcommand\v[2][\empty]{\vm{#2}_{#1}}
\newcommand\dv[2][\empty]{\dot{\vm{#2}}_{#1}}
\newcommand\ddv[2][\empty]{\ddot{\vm{#2}}_{#1}}

% Scalars with optional index and derivatives
\newcommand\s[2][\empty]{#2_{#1}}
\newcommand\ds[2][\empty]{\dot{#2}_{#1}}
\newcommand\dds[2][\empty]{\ddot{#2}_{#1}}

% Vectors tilde with optional index and derivatives

\newcommand\vt[2][\empty]{\tilde{\vm{#2}}_{#1}}
\newcommand\dvt[2][\empty]{\dot{\tilde{\vm{#2}}}_{#1}}
\newcommand\ddvt[2][\empty]{\ddot{\tilde{\vm{#2}}}_{#1}}

% Vectors hat with optional index and derivatives
\newcommand\vh[2][\empty]{\hat{\vm{#2}}_{#1}}
\newcommand\dvh[2][\empty]{\dot{\hat{\vm{#2}}}_{#1}}
\newcommand\ddvh[2][\empty]{\ddot{\hat{\vm{#2}}}_{#1}}

% Vectors bar with optional index and derivatives
\newcommand\vb[2][\empty]{\bar{\vm{#2}}_{#1}}
\newcommand\dvb[2][\empty]{\dot{\bar{\vm{#2}}}_{#1}}
\newcommand\ddvb[2][\empty]{\ddot{\bar{\vm{#2}}}_{#1}}

% Scalars tilde with optional index and derivatives
\newcommand\st[2][\empty]{\tilde{#2}_{#1}}
\newcommand\dst[2][\empty]{\dot{\tilde{#2}}_{#1}}
\newcommand\ddst[2][\empty]{\ddot{\tilde{#2}}_{#1}}

% Scalars hat with optional index and derivatives
\newcommand\sh[2][\empty]{\hat{#2}_{#1}}
\newcommand\dsh[2][\empty]{\dot{\hat{#2}}_{#1}}
\newcommand\ddsh[2][\empty]{\ddot{\hat{#2}}_{#1}}

% Scalars bar with optional index and derivatives
% note: the first of these commands overwrite the "\sb{}" command which can be used instead of _{} to yield subscript. But who uses that?
\renewcommand\sb[2][\empty]{\bar{#2}_{#1}}
\newcommand\dsb[2][\empty]{\dot{\bar{#2}}_{#1}}
\newcommand\ddsb[2][\empty]{\ddot{\bar{#2}}_{#1}}

\newcommand\conj[1]{\overline{#1}}
\newcommand\upto[2][0]{\set{#1\dots #2}}
\newcommand\ceil[1]{\left\lceil#1\right\rceil}
\newcommand\floor[1]{\left\lfloor#1\right\rfloor}

%maybe figure these two out
\newcommand\der[2][]{\frac{\mathrm{d} #1}{\mathrm{d} #2}}
\newcommand\pder[2][]{\frac{\partial#1}{\partial#2}}
\newcommand\lder{\mathcal{L}}

\let\atan\arctan
\let\vphi\varphi
\DeclareMathOperator{\diag}{diag}

% Useful math operators
\DeclareMathOperator{\Ad}{Ad}
\DeclareMathOperator{\ad}{ad}
\DeclareMathOperator{\sgn}{\mathrm{sgn}}
\DeclareMathOperator{\blkdiag}{blkdiag}
\DeclareMathOperator{\SE}{SE}
\DeclareMathOperator{\se}{\mathfrak{se}}
\DeclareMathOperator{\SO}{SO}
\DeclareMathOperator{\so}{\mathfrak{so}}
\DeclareMathOperator{\rank}{rank}


\newcommand{\norm}[1]{\left\lVert#1\right\rVert}
\newcommand{\abs}[1]{\left|#1\right|}
\newcommand{\absign}[1]{\lceil#1\rfloor} % Signed absolute value



\gdef\i{\ensuremath{\mathrm{i}}}

% \DeclareMathOperator{\uh}{^{\wedge}}
\def\uh{^{\wedge}}
% \DeclareMathOperator{\lh}{^{\vee}}
\def\lh{^{\vee}}

% \newcommand\abs[1]{\left|#1\right|}

\makeatletter
\renewenvironment{proof}[1][\relax]{\par
  \pushQED{\qed}%
  \normalfont \topsep6\p@\@plus6\p@\relax
  \trivlist
  \item[\hskip\labelsep\bfseries
    \ifx#1\relax \proofname\else#1\fi\@addpunct{.}]\ignorespaces
}{%
  \popQED\endtrivlist\@endpefalse
}
\makeatother

\renewcommand{\qedsymbol}{$\blacksquare$}

\newtheorem{lemma}{Lemma}
\newtheorem{theorem}{Theorem}
\newtheorem{remark}{Remark}
\newtheorem{assumption}{Assumption}
\newtheorem{definition}{Definition}
\newtheorem{corollary}{Corollary}
\newtheorem{proposition}{Proposition}
%\newtheorem{prop}{Proposition}


\makeatletter
\providecommand{\examplesymbol}{\ensuremath\blacktriangle}
\newenvironment{example}[1][\examplename]{\par
  \pushQED{\qed}%
  \normalfont \topsep6\p@\@plus6\p@\relax
  \begingroup
  \let\qedsymbol\examplesymbol
      \trivlist

  \item[\hskip\labelsep
        \itshape
    #1\@addpunct{.}]\ignorespaces
}{%
  \popQED\endtrivlist\@endpefalse
  \endgroup
}
\providecommand{\examplename}{Example}
\makeatother

\crefname{prop}{Proposition}{Propositions}
\crefname{lemma}{Lemma}{Lemmas}
\crefname{theorem}{Theorem}{Theorems}
\crefname{assumption}{Assumption}{Assumptions}
\crefname{definition}{Definition}{Definitions}
\crefname{remark}{Remark}{Remarks}
\crefname{equation}{}{}
\crefname{corollary}{Corollary}{Corollaries}


% Adjusts letter placement etc, seems to fix a lot of over/underfull issues
\usepackage[final]{microtype}